% Part: first-order-logic
% Chapter: proof-systems
% Section: seqeunt-calculus

\documentclass[../../../include/open-logic-section]{subfiles}

\begin{document}

\iftag{FOL}
      {\olfileid{fol}{prf}{seq}}
      {\olfileid{pl}{prf}{seq}}

\olsection{حساب التتابعيات}

في حين تعمل أنظمة !!{derivation} كثيرة بترتيبات من ال!!{sentence}s،
يعمل حساب التتابعيات بـ\emph{التتابعيات}. والتتابعية تعبير من الصورة
\[
!A_1, \dots, !A_m \Sequent !B_1, \dots, !B_m,
\]
أي إنها زوج من متتاليتين من ال!!{sentence}s يفصل بينهما رمز
التتابعية~$\Sequent$. ويجوز أن تكون أي من المتتاليتين خالية.
و!!^a{derivation} في حساب التتابعيات شجرة من التتابعيات؛ تكون
التتابعيات العليا فيها من صورة خاصة (وتسمى «التتابعيات الابتدائية» أو
«البديهيات»)، وتتبع كل تتابعية أخرى من التتابعيات الواقعة فوقها مباشرةً
بإحدى قواعد الاستدلال. فإما أن تعالج قواعد الاستدلال ال!!{sentence}s
في التتابعيات (بإضافتها أو حذفها أو إعادة ترتيبها في الجانب الأيسر أو
الأيمن)، وإما أن تُدخل !!{formula} مركبة في نتيجة القاعدة. فمثلًا،
تجيز القاعدة~$\LeftR{\land}$ الاستدلال من $!A,
\Gamma \Sequent \Delta$ إلى $!A \land !B, \Gamma \Sequent \Delta$،
وتجيز القاعدة~$\RightR{\lif}$ الاستدلال من $!A, \Gamma \Sequent
\Delta, !B$ إلى $\Gamma \Sequent \Delta, !A \lif !B$، لأي من
$\Gamma$ و$\Delta$ و$!A$ و~$!B$. (ويجوز خصوصًا أن تكون $\Gamma$
و~$\Delta$ خاليتين.)

تُعرّف علاقة~$\Proves$ القائمة على حساب التتابعيات كما يأتي:
$\Gamma \Proves !A$ إذا وفقط إذا وجدت متتالية~$\Gamma_0$ بحيث يكون
كل~$!A$ في~$\Gamma_0$ عنصرًا في~$\Gamma$، ويوجد !!{derivation} تقع
التتابعية~$\Gamma_0 \Sequent !A$ في جذره. وتكون~$!A$ مبرهنة في حساب
التتابعيات إذا كان للتتابعية~$\Sequent !A$ !!a{derivation}. فمثلًا،
هذا !!a{derivation} يبين أن $\Proves (!A \land !B) \lif !A$:
\begin{prooftree}
  \Axiom$!A \fCenter !A$
  \RightLabel{\LeftR{\land}}
  \UnaryInf$!A \land !B \fCenter !A$
  \RightLabel{\RightR{\lif}}
  \UnaryInf$\fCenter (!A \land !B) \lif !A$
\end{prooftree}

تكون المجموعة~$\Gamma$ غير متسقة في حساب التتابعيات إذا وجد
!!a{derivation} لـ~$\Gamma_0 \Sequent$ (حيث يكون كل
$!A \in \Gamma_0$ عنصرًا في~$\Gamma$، ويكون الجانب الأيمن من
التتابعية خاليًا). وباستعمال القاعدة~\RightR{\Weakening} تصبح أي
!!{sentence} قابلةً لل!!{derive} من مجموعة غير متسقة.

ابتكر غيرهارد غنتسن حساب التتابعيات في ثلاثينيات القرن العشرين. وبفضل
تصميمه المنهجي والمتناظر، فإنه جهاز صوري نافع جدًا لتطوير نظرية
لل!!{derivation}s. ويسهل نسبيًا العثور على !!{derivation}s في حساب
التتابعيات، لكن قراءة هذه ال!!{derivation}s كثيرًا ما تكون عسيرة، ولا يسهل أحيانًا رؤية
صلتها بالبراهين. ومع ذلك، فقد ثبت أنه مقاربة أنيقة جدًا لأنظمة
ال!!{derivation}، ولكثير من الأنساق المنطقية أنظمة من حساب
التتابعيات.

\end{document}
